\plannedchapter{栈、堆、运行时和对象生命周期}
{本章把函数调用栈、局部变量、堆分配、全局构造、RAII 和运行时支持连接起来，为 ABI 和性能测量打基础。}
{\item 函数调用为什么需要栈？
\item 栈帧里可能有哪些内容？
\item \code{new}/\code{malloc} 为什么不能随便放在 hot loop？}
{\item call/ret、返回地址、保存寄存器、局部变量和栈对齐。
\item allocator、free list、系统调用和线程缓存的粗略模型。
\item 全局构造、析构、异常展开和运行时库。}
{用 GDB 观察 \code{rsp/rbp}；写递归和非递归函数比较栈增长；benchmark 循环内分配与预分配。}
